% Источники: exam/materials/моделирование.pdf, раздел 4; exam/materials/преподаватель/2020/01-05-2020-Model_L11__1_.pdf; exam/materials/преподаватель/2020/01-05-2020-Model_L11__2_.pdf; exam/materials/прошлогодние ответы/Modelirovanie_1_-3.pdf, вопрос 34; GitHub HumsterProgrammer/iu9-notes/8sem/Моделирование/Моделирование_лекция-15_02-04-26.md.
\section{Планирование эксперимента. История развития научного направления.}

\textbf{Планирование эксперимента} — комплекс мероприятий, направленных на эффективное проведение эксперимента, в том числе вычислительного.

Планирование эксперимента заключается в поиске оптимальных условий, построении интерполяционных формул, выборе значимых факторов, оценке и уточнении констант теоретических моделей.

Основная цель планирования эксперимента — достижение максимальной точности измерений при минимальном количестве проведенных опытов и сохранении статистической достоверности результатов.

Решение задач теории планирования эксперимента предусматривает использование априорной информации об изучаемом процессе для выбора общей последовательности управления экспериментами, которая уточняется после очередного этапа проведения исследований на основе вновь полученных сведений. Тем самым достигается возможность рационального управления экспериментами при неполном первоначальном знании характеристик исследуемого объекта.

Планирование эксперимента возникло как научное направление в 20-х гг. XX века на основе работ Р.~Фишера (\emph{Design of Experiments}). Ряд исследователей: Кишен, Манн, Бойс, Гёлс, Рао и другие — обосновали практические достижения Фишера и оптимизировали их.

С конца 40-х — начала 50-х гг. планирование эксперимента активно развивалось в США: Бокс, Уилсон, Юл и другие революционно упростили подход к решению задачи планирования эксперимента, сводя их к степенным разложениям и методу наименьших квадратов.

Развитие теории планирования эксперимента получило развитие в категории ядерной техники (Соколов, Федоров). После чего распространилось во всех сферах науки и промышленности. Тогда же появились монографии по планированию эксперимента (Адлер, Вороновский, Маркова).

\clearpage
